% !TEX program = pdflatex
% Источник: курс Hugging Face Deep Reinforcement Learning, модуль 1.
% Compile from the notes/ directory (default LaTeX Workshop + latexmk -pdf):
%   cd notes/ru
%   latexmk -pdf -interaction=nonstopmode unit1ru.tex
% Файлы ../preamble.tex и ../assets/unit1/ доступны из notes/ru/.
\PassOptionsToPackage{T2A}{fontenc}
\PassOptionsToPackage{russian}{babel}
\documentclass[11pt,a4paper]{article}

% Shared typography/layout for Hugging Face Deep RL Course notes.
% Compile with pdfLaTeX / latexmk -pdf with notes/ as the working directory.
% e.g. cd notes && latexmk -pdf unit1.tex

\usepackage[T1]{fontenc}
\usepackage[utf8]{inputenc}
\usepackage[english]{babel}
\usepackage{lmodern}
\usepackage{textcomp}
\usepackage{geometry}
\geometry{a4paper,margin=22mm,headheight=14pt,headsep=8mm,footskip=12mm}

\usepackage{amsmath,amssymb,mathtools}
\usepackage{graphicx,caption,needspace,placeins}
\usepackage{array,booktabs,tabularx}
\usepackage{xcolor}
\usepackage{microtype}
\usepackage{enumitem}
\usepackage{listings}
\usepackage[most]{tcolorbox}
\usepackage{titlesec}
\usepackage{fancyhdr}
\usepackage{hyperref}
\usepackage{bookmark}

\definecolor{ink}{HTML}{243041}
\definecolor{accent}{HTML}{38598C}
\definecolor{secondary}{HTML}{60768A}
\definecolor{codebg}{HTML}{F4F6F8}
\definecolor{codeborder}{HTML}{E3E8EE}
\definecolor{notebg}{HTML}{F3F7FC}

\hypersetup{
  colorlinks=true,
  linkcolor=accent,
  urlcolor=accent,
  citecolor=accent,
  pdfborder={0 0 0}
}
\urlstyle{same}

\setlength{\parindent}{0pt}
\setlength{\parskip}{0.38em plus 0.12em minus 0.05em}

\setlist[itemize]{
  topsep=0.35em,
  itemsep=0.13em,
  leftmargin=1.7em
}
\setlist[enumerate]{
  topsep=0.35em,
  itemsep=0.13em,
  leftmargin=1.8em
}

\setcounter{tocdepth}{2}
\setcounter{secnumdepth}{2}

\titleformat{\section}
  {\Large\bfseries\color{ink}}
  {\thesection}{0.7em}{}

\titleformat{\subsection}
  {\large\bfseries\color{accent}}
  {\thesubsection}{0.7em}{}

\titleformat{\subsubsection}
  {\normalsize\bfseries\color{ink}}
  {\thesubsubsection}{0.6em}{}

\titlespacing*{\section}
  {0pt}{2.2ex plus 0.7ex}{0.9ex}
\titlespacing*{\subsection}
  {0pt}{1.7ex plus 0.4ex}{0.6ex}
\titlespacing*{\subsubsection}
  {0pt}{1.3ex plus 0.3ex}{0.45ex}

\pagestyle{fancy}
\fancyhf{}

\fancyhead[L]{
  \footnotesize\sffamily\color{secondary}
  HF Deep RL Course \enspace/\enspace Unit 1
}
\fancyhead[R]{
  \footnotesize\sffamily\color{secondary}
  \nouppercase{\rightmark}
}
\fancyfoot[C]{
  \footnotesize\color{secondary}\thepage
}
\renewcommand{\headrulewidth}{0.25pt}

\captionsetup{
  font=small,
  labelfont={bf,color=secondary},
  textfont={color=ink},
  justification=centering,
  skip=4pt,
  hypcap=false
}

\lstdefinestyle{hfcode}{
  basicstyle=\ttfamily\footnotesize,
  backgroundcolor=\color{codebg},
  frame=single,
  rulecolor=\color{codeborder},
  numbers=none,
  showstringspaces=false,
  breaklines=true,
  breakatwhitespace=false,
  columns=fullflexible,
  keepspaces=true,
  xleftmargin=0pt,
  xrightmargin=0pt,
  aboveskip=0.6em,
  belowskip=0.6em,
  framexleftmargin=7pt,
  framexrightmargin=7pt,
  framextopmargin=5pt,
  framexbottommargin=5pt,
  upquote=true,
  tabsize=4,
}
\lstset{style=hfcode}

\newcommand{\CourseFigure}[2]{%
  \par\addvspace{0.55em}\noindent%
  \begin{minipage}{\linewidth}%
    \centering
    \includegraphics[
      width=0.86\linewidth,
      height=0.295\textheight,
      keepaspectratio
    ]{#1}\par
    \captionof{figure}{#2}%
  \end{minipage}%
  \par\addvspace{0.55em}%
}

\newtcolorbox{coursenote}{
  enhanced,
  breakable,
  colback=notebg,
  colframe=accent!30,
  boxrule=0.35pt,
  arc=1.2mm,
  left=9pt,
  right=9pt,
  top=6pt,
  bottom=6pt
}

\renewenvironment{quote}
  {\begin{coursenote}}
  {\end{coursenote}}

\emergencystretch=2.5em

\providecommand{\tightlist}{
  \setlength{\itemsep}{0pt}
  \setlength{\parskip}{0pt}
}

\providecommand{\passthrough}[1]{#1}

\Urlmuskip=0mu plus 6mu\relax

\renewcommand{\rmdefault}{cmr}
\renewcommand{\sfdefault}{cmss}
\renewcommand{\ttdefault}{cmtt}
\graphicspath{{../assets/unit1/}}
\fancyhead[L]{\footnotesize\sffamily\color{secondary}Курс Deep RL\enspace/\enspace Модуль 1}
\title{\textbf{Глубокое обучение с подкреплением}\\[0.35em]\large Модуль 1: основы и первый агент}
\author{Курс Deep RL от Hugging Face\\\small Полное учебное издание на основе исходных MDX-файлов}
\date{}
\begin{document}
\selectlanguage{russian}
\fontencoding{T2A}\selectfont
\maketitle
\vspace{-0.8em}
\begin{center}\small\color{secondary}Исходный материал: курс Deep RL от Hugging Face.\end{center}
\tableofcontents
\vspace{1em}
\clearpage
\section{Введение в глубокое обучение с подкреплением}\label{introduction-to-deep-reinforcement-learning}

\CourseFigure{thumbnail.jpg}{Модуль 1: введение в глубокое обучение с подкреплением (официальная иллюстрация курса)}

Добро пожаловать в одну из самых увлекательных областей искусственного интеллекта: \textbf{глубокое обучение с подкреплением.}

Глубокое обучение с подкреплением (Deep RL) --- это разновидность машинного обучения, при которой агент учится \textbf{правильно действовать} в среде, \textbf{совершая действия} и \textbf{наблюдая их результаты.}

В этом первом модуле \textbf{вы изучите основы глубокого обучения с подкреплением.}

Затем с помощью \href{https://stable-baselines3.readthedocs.io/en/master/}{Stable-Baselines3}, библиотеки для глубокого обучения с подкреплением, вы \textbf{обучите агента, управляющего лунным посадочным модулем, правильно приземляться на Луну.}

\CourseFigure{lunarLander-still.jpg}{Среда LunarLander: статичный кадр с официальной обложки курса; оригинальная анимация доступна по \href{https://huggingface.co/datasets/huggingface-deep-rl-course/course-images/resolve/main/en/unit1/lunarLander.gif}{ссылке на GIF}}

И наконец, вы \textbf{загрузите обученного агента в Hugging Face Hub --- бесплатную открытую платформу, где можно делиться моделями машинного обучения, наборами данных и демонстрациями.}

Прежде чем приступать к реализации агентов глубокого обучения с подкреплением, очень важно \textbf{освоить эти понятия}. Цель этой главы --- дать вам прочную теоретическую основу.

После завершения модуля в дополнительном разделе вы также \textbf{сможете обучить пса Huggy приносить палку и поиграть с ним.}

\CourseFigure{huggy-bonus.jpg}{Пёс Huggy: кадр из официальных материалов курса (\href{https://huggingface.co/datasets/huggingface-deep-rl-course/course-images/resolve/main/en/unit0/huggy.mp4}{оригинальное видео})}

Приступим!

\section{Что такое обучение с подкреплением?}\label{what-is-reinforcement-learning}

Чтобы разобраться в обучении с подкреплением, сначала посмотрим на общую картину.

\subsection{Общая идея}\label{the-big-picture}

Основная идея обучения с подкреплением заключается в том, что агент (система искусственного интеллекта) учится в среде, \textbf{взаимодействуя с ней} методом проб и ошибок и \textbf{получая вознаграждения} (положительные или отрицательные) в качестве обратной связи за совершённые действия.

Обучение через взаимодействие со средой \textbf{знакомо нам по повседневному опыту.}

Представьте, например, что вы посадили младшего брата перед видеоигрой, в которую он никогда не играл, дали ему геймпад и оставили одного.

\CourseFigure{Illustration_1.jpg}{Первое действие: игрок движется вправо}

Ваш брат взаимодействует со средой (видеоигрой), нажимая кнопку движения вправо (действие). Он подбирает монету и получает вознаграждение +1. Это положительный результат: он только что понял, что в этой игре \textbf{нужно собирать монеты.}

\CourseFigure{Illustration_2.jpg}{Положительное вознаграждение: подбор монеты}

Но затем \textbf{он снова нажимает кнопку вправо} и сталкивается с противником. Персонаж погибает, поэтому вознаграждение составляет -1.

\CourseFigure{Illustration_3.jpg}{Отрицательное вознаграждение: столкновение с противником}

Взаимодействуя со средой методом проб и ошибок, ваш младший брат понимает, что \textbf{нужно собирать монеты, но избегать врагов.}

\textbf{Без какого-либо внешнего руководства} ребёнок постепенно начинает играть всё лучше.

Так же учатся люди и животные --- \textbf{через взаимодействие.} Обучение с подкреплением --- это \textbf{вычислительный подход к обучению на основе действий.}

\subsubsection{Формальное определение}\label{a-formal-definition}

Теперь можно дать формальное определение:

\begin{quote}
\textbf{Примечание к курсу.} Обучение с подкреплением --- это подход к решению задач управления (также называемых задачами принятия решений), в котором создаются агенты, обучающиеся при взаимодействии со средой методом проб и ошибок и получающие в качестве единственной обратной связи вознаграждения --- положительные или отрицательные.
\end{quote}

Но как именно работает обучение с подкреплением?

\section{Формальная схема обучения с подкреплением}\label{the-reinforcement-learning-framework}

\subsection{Процесс RL}\label{the-rl-process}

\CourseFigure{RL_process.jpg}{Цикл RL: состояние, действие, вознаграждение и следующее состояние. Источник: Reinforcement Learning: An Introduction, Richard Sutton и Andrew G. Barto}

\noindent{\footnotesize Источник: \url{http://incompleteideas.net/book/RLbook2020.pdf}}

Чтобы понять процесс RL, представим агента, который учится играть в платформер:

\CourseFigure{RL_process_game.jpg}{Цикл взаимодействия в RL на примере видеоигры}

\begin{itemize}
\tightlist
\item
  Наш агент получает от \textbf{среды} \textbf{состояние \(S_0\)} --- первый кадр игры.
\item
  На основе этого \textbf{состояния \(S_0\)} агент выбирает \textbf{действие \(A_0\)} --- движется вправо.
\item
  Среда переходит в \textbf{новое состояние \(S_1\)} --- следующий кадр.
\item
  Среда выдаёт агенту некоторое \textbf{вознаграждение \(R_1\)} --- мы не погибли, то есть получили \emph{положительное вознаграждение +1}.
\end{itemize}

Такой цикл RL создаёт последовательность \textbf{состояний, действий, вознаграждений и следующих состояний.}

\CourseFigure{sars.jpg}{Кортеж перехода: состояние, действие, вознаграждение, следующее состояние}

Цель агента --- \emph{максимизировать} суммарное вознаграждение, \textbf{называемое ожидаемой отдачей (expected return).}

\subsection{Гипотеза вознаграждения: центральная идея RL}\label{the-reward-hypothesis-the-central-idea-of-reinforcement-learning}

\(\Rightarrow\) Почему агент должен максимизировать ожидаемую отдачу?

Потому что RL основывается на \textbf{гипотезе вознаграждения}: любую цель можно описать как \textbf{максимизацию ожидаемой отдачи} (ожидаемого суммарного вознаграждения).

Поэтому, чтобы \textbf{получить наилучшее поведение}, в RL мы учимся выбирать действия, \textbf{максимизирующие ожидаемое суммарное вознаграждение.}

\subsection{Марковское свойство}\label{markov-property}

В научных работах процесс RL часто описывается как \textbf{марковский процесс принятия решений} (Markov Decision Process, MDP).

К марковскому свойству мы ещё вернёмся в следующих модулях. Пока достаточно запомнить: благодаря ему для \textbf{выбора действия агенту нужно знать только текущее состояние}, а \textbf{не всю историю предыдущих состояний и действий.}

\subsection{Пространство наблюдений и состояний}\label{observationsstates-space}

Наблюдения и состояния --- это \textbf{информация, которую агент получает от среды.} Для видеоигры такой информацией может быть кадр (скриншот), а для торгового агента --- цена определённой акции.

Однако между \emph{наблюдением} и \emph{состоянием} существует различие:

\begin{itemize}
\tightlist
\item
  \emph{Состояние \(s\)} --- это \textbf{полное описание мира} (без скрытой информации). Оно доступно в полностью наблюдаемой среде.
\end{itemize}

\CourseFigure{chess.jpg}{В шахматах мы получаем от среды состояние, поскольку имеем доступ ко всей информации на доске}

В шахматах нам известна вся расстановка фигур, поэтому среда предоставляет состояние. Иными словами, она полностью наблюдаема.

\begin{itemize}
\tightlist
\item
  \emph{Наблюдение \(o\)} --- это \textbf{частичное описание состояния}. Оно доступно в частично наблюдаемой среде.
\end{itemize}

\CourseFigure{mario.jpg}{В Super Mario Bros видна только часть уровня рядом с персонажем, поэтому агент получает наблюдение}

В Super Mario Bros мы видим только ближайший к персонажу участок уровня, поэтому получаем наблюдение.

Таким образом, Super Mario Bros --- частично наблюдаемая среда: мы получаем наблюдение, \textbf{поскольку видим только часть уровня.}

\begin{quote}
\textbf{Примечание к курсу.} В этом курсе слово «состояние» используется и для полного состояния, и для наблюдения, однако при реализации алгоритмов мы будем различать эти понятия.
\end{quote}

Подведём итог:

\CourseFigure{obs_space_recap.jpg}{Полное состояние и частичное наблюдение}

\subsection{Пространство действий}\label{action-space}

Пространство действий --- множество \textbf{всех возможных действий в данной среде.}

Действия могут принадлежать \emph{дискретному} или \emph{непрерывному пространству}:

\begin{itemize}
\tightlist
\item
  \emph{Дискретное пространство}: число возможных действий \textbf{конечно}.
\end{itemize}

\CourseFigure{mario.jpg}{В Super Mario Bros всего четыре возможных действия: влево, вправо, вверх (прыжок) и вниз (приседание)}

В Super Mario Bros множество действий конечно, поскольку доступны только четыре направления.

\begin{itemize}
\tightlist
\item
  \emph{Непрерывное пространство}: число возможных действий \textbf{бесконечно}.
\end{itemize}

\CourseFigure{self_driving_car.jpg}{У агента, управляющего автомобилем, бесконечное множество возможных действий: повернуть на 20\textdegree{}, 21,1\textdegree{}, 21,2\textdegree{}, подать звуковой сигнал, повернуть вправо на 20\textdegree{} и т. д.}

Подведём итог:

\CourseFigure{action_space.jpg}{Дискретные и непрерывные пространства действий}

Эту информацию важно учитывать, поскольку она \textbf{влияет на выбор алгоритма RL в дальнейшем.}

\subsection{Вознаграждения и дисконтирование}\label{rewards-and-the-discounting}

Вознаграждение играет фундаментальную роль в RL, поскольку это \textbf{единственная обратная связь} для агента. Благодаря ему агент понимает, \textbf{насколько удачным было совершённое действие.}

Суммарное вознаграждение для момента времени \textbf{t} можно записать так:

\CourseFigure{rewards_1.jpg}{Суммарное вознаграждение равно сумме всех вознаграждений в последовательности}

Что эквивалентно выражению:

\CourseFigure{rewards_2.jpg}{Суммарное вознаграждение: \(r_{t+1}+r_{t+2}+\cdots\), где \(r_{t+k+1}\) --- вознаграждение на соответствующем шаге}

Однако в реальности \textbf{нельзя просто складывать все вознаграждения таким образом.} Более ранние вознаграждения (в начале игры) \textbf{с большей вероятностью будут получены}, поскольку они предсказуемее отдалённых будущих вознаграждений.

Представим, что агент --- маленькая мышь, которая за один шаг может переместиться на одну клетку, а её противник --- кот (который тоже движется). Цель мыши --- \textbf{съесть как можно больше сыра до того, как её поймает кот.}

\CourseFigure{rewards_3.jpg}{Мышь, сыр и кот: иллюстрация отложенных вознаграждений}

Как видно на схеме, \textbf{съесть сыр поблизости вероятнее, чем сыр рядом с котом} (чем ближе мышь к коту, тем опаснее).

Поэтому \textbf{вознаграждение рядом с котом, даже если оно больше (сыра там больше), будет сильнее дисконтироваться}: мы не уверены, что сможем до него добраться.

Для дисконтирования вознаграждений действуем следующим образом:

\begin{enumerate}
\def\labelenumi{\arabic{enumi}.}
\tightlist
\item
  Вводим коэффициент дисконтирования \(\gamma\). Он \textbf{должен лежать между 0 и 1}; чаще всего используют значения \textbf{от 0,95 до 0,99}.
\end{enumerate}

\begin{itemize}
\tightlist
\item
  Чем больше \(\gamma\), тем слабее дисконтирование. Это означает, что агент \textbf{сильнее учитывает долгосрочное вознаграждение.}
\item
  Напротив, чем меньше \(\gamma\), тем сильнее дисконтирование, то есть агент \textbf{больше ориентируется на ближайшее вознаграждение (ближний сыр).}
\end{itemize}

\begin{enumerate}
\def\labelenumi{\arabic{enumi}.}
\setcounter{enumi}{1}
\tightlist
\item
  Затем каждое вознаграждение умножается на \(\gamma\), возведённую в степень, соответствующую номеру шага. Чем дальше шаг во времени, тем ближе к мыши может оказаться кот, \textbf{поэтому вероятность получить будущее вознаграждение всё ниже.}
\end{enumerate}

Наша ожидаемая дисконтированная суммарная отдача имеет вид:

\CourseFigure{rewards_4.jpg}{Дисконтированная отдача с коэффициентом дисконтирования \(\gamma\)}

\section{Типы задач}\label{type-of-tasks}

Задача --- это \textbf{конкретный экземпляр} проблемы RL. Различают два типа задач: \textbf{эпизодические} и \textbf{непрерывные}.

\subsection{Эпизодические задачи}\label{episodic-task}

В этом случае есть начальная и конечная точки \textbf{(терминальное состояние). Вместе они задают эпизод}: последовательность состояний, действий, вознаграждений и новых состояний.

Например, в Super Mario Bros эпизод начинается при запуске уровня и заканчивается, \textbf{когда персонаж погибает или достигает конца уровня.}

\CourseFigure{mario.jpg}{Начало нового эпизода}

\subsection{Продолжающиеся задачи}\label{continuing-tasks}

Это задачи, которые продолжаются бесконечно (\textbf{без терминального состояния}). Здесь агент должен \textbf{научиться выбирать лучшие действия, одновременно продолжая взаимодействовать со средой.}

Например, агент автоматической торговли акциями: у такой задачи нет заранее заданного конца. \textbf{Агент работает, пока мы сами не остановим его.}

\CourseFigure{stock.jpg}{Биржевая торговля как пример продолжающейся задачи}

Подведём итог:

\CourseFigure{tasks.jpg}{Эпизодические и продолжающиеся задачи}

\section{Компромисс между исследованием и использованием знаний}\label{the-explorationexploitation-trade-off}

Прежде чем рассматривать способы решения задач RL, необходимо обсудить ещё одну очень важную тему: \emph{компромисс между исследованием (exploration) и использованием уже полученных знаний (exploitation).}

\begin{itemize}
\tightlist
\item
  \emph{Исследование} --- изучение среды с помощью случайных действий для \textbf{получения дополнительной информации о ней.}
\item
  \emph{Использование знаний} --- \textbf{применение уже известной информации для максимизации вознаграждения.}
\end{itemize}

Помните, что цель агента --- максимизировать ожидаемое суммарное вознаграждение. Однако \textbf{здесь легко попасть в распространённую ловушку}.

Рассмотрим пример:

\CourseFigure{exp_1.jpg}{Сырный лабиринт: локальные и отдалённые вознаграждения при выборе между исследованием и использованием знаний}

В этой игре мышь может получать \textbf{неограниченное количество маленьких кусочков сыра} (по +1 каждый). Но в верхней части лабиринта находится огромная куча сыра (+1000).

Однако если заниматься только использованием известных возможностей, агент никогда не доберётся до огромной кучи. Он будет получать \textbf{ближайшие вознаграждения}, даже если они совсем небольшие (exploitation).

Если же агент немного исследует среду, он сможет \textbf{обнаружить большое вознаграждение} (огромную кучу сыра).

Именно это и называют компромиссом exploration/exploitation. Нужно соблюдать баланс между тем, насколько активно мы \textbf{исследуем среду} и насколько \textbf{используем уже известную информацию о ней.}

Поэтому необходимо \textbf{задать правило управления этим компромиссом}. В следующих модулях мы рассмотрим разные способы его реализации.

Если идея пока неочевидна, \textbf{вспомните повседневную задачу --- выбор ресторана}:

\CourseFigure{exp_2.jpg}{Исследование и использование знаний: знакомый ресторан или неизвестная альтернатива. Источник: курс Berkeley AI}

\noindent{\footnotesize Источник: \url{https://inst.eecs.berkeley.edu/~cs188/sp20/assets/lecture/lec15_6up.pdf}}

\begin{itemize}
\tightlist
\item
  \emph{Использование знаний}: каждый день ходить в уже проверенное хорошее место, \textbf{рискуя не найти другой ресторан, который мог бы понравиться ещё больше.}
\item
  \emph{Исследование}: пробовать новые рестораны, рискуя получить плохой опыт, \textbf{но сохраняя шанс открыть отличное место.}
\end{itemize}

Подведём итог:

\CourseFigure{expexpltradeoff.jpg}{Исследование и использование знаний: обобщающая схема}

\section{Два основных подхода к решению задач RL}\label{two-main-approaches-for-solving-rl-problems}

\begin{quote}
\textbf{Примечание к курсу.} Теперь, когда мы познакомились с основами RL, как же решать соответствующие задачи?
\end{quote}

Другими словами, как построить агента, способного \textbf{выбирать действия, максимизирующие его ожидаемое суммарное вознаграждение?}

\subsection{Политика \texorpdfstring{$\pi$}{pi}: «мозг» агента}\label{the-policy-ux3c0-the-agents-brain}

Политика \textbf{$\pi$} --- это \textbf{«мозг» агента}, функция, определяющая, \textbf{какое действие нужно выбрать в данном состоянии.} Таким образом, она \textbf{задаёт поведение агента} в каждый момент времени.

\CourseFigure{policy_1.jpg}{Политика как «мозг» агента: функция, определяющая действие по текущему состоянию}

Именно эту \textbf{функцию политики нам необходимо выучить}. Наша цель --- найти оптимальную политику $\pi^*$, которая \textbf{максимизирует ожидаемую отдачу}, если агент ей следует. Оптимальную политику $\pi^*$ мы находим \textbf{в процессе обучения.}

Существует два подхода к обучению агента для поиска оптимальной политики $\pi^*$:

\begin{itemize}
\tightlist
\item
  \textbf{Непосредственно} обучать агента выбирать \textbf{действие по текущему состоянию}: \textbf{методы на основе политики (Policy-Based Methods).}
\item
  Косвенно: \textbf{обучать агента оценивать ценность состояний}, чтобы затем выбирать действия, \textbf{ведущие к более ценным состояниям}: методы на основе функции ценности (Value-Based Methods).
\end{itemize}

\subsection{Методы на основе политики}\label{policy-based-methods}

В методах на основе политики \textbf{мы непосредственно обучаем функцию политики.}

Эта функция сопоставляет каждому состоянию наиболее подходящее действие. Другой вариант --- определять \textbf{распределение вероятностей по возможным действиям в данном состоянии.}

\CourseFigure{policy_2.jpg}{Детерминированная политика непосредственно указывает действие на каждом шаге}

Различают два вида политик:

\begin{itemize}
\tightlist
\item
  \emph{Детерминированная}: в заданном состоянии политика \textbf{всегда возвращает одно и то же действие.}
\end{itemize}

\CourseFigure{policy_3.jpg}{Действие = политика(состояние)}

\CourseFigure{policy_4.jpg}{Стохастическая политика: действия выбираются из условного распределения вероятностей}

\begin{itemize}
\tightlist
\item
  \emph{Стохастическая}: возвращает \textbf{распределение вероятностей по действиям.}
\end{itemize}

\CourseFigure{policy_5.jpg}{Политика(действия | состояние) = распределение вероятностей возможных действий при текущем состоянии}

\CourseFigure{policy-based.png}{Для заданного начального состояния стохастическая политика возвращает распределение вероятностей всех доступных действий}

Подведём итог:

\CourseFigure{pbm_1.jpg}{Методы на основе политики: детерминированная политика}

\CourseFigure{pbm_2.jpg}{Методы на основе политики: стохастическая политика}

\subsection{Методы на основе функции ценности}\label{value-based-methods}

В методах на основе функции ценности вместо непосредственного обучения политики \textbf{мы обучаем функцию ценности}, которая сопоставляет состоянию ожидаемую \textbf{ценность нахождения в нём.}

Ценность состояния --- это \textbf{ожидаемая дисконтированная отдача}, которую агент может получить, \textbf{начав из этого состояния и в дальнейшем следуя нашей политике.}

Фраза «следуя нашей политике» здесь означает, что она будет \textbf{«выбирать состояние с наибольшей ценностью».}

\CourseFigure{value_1.jpg}{Функция ценности состояния и ожидаемая дисконтированная отдача}

На рисунке видно, что наша функция ценности \textbf{назначила значения всем возможным состояниям.}

\CourseFigure{value_2.jpg}{Благодаря функции ценности политика на каждом шаге выбирает состояние с наибольшей оценкой: -7, затем -6, затем -5 и так далее, пока не достигнет цели}

Благодаря функции ценности политика на каждом шаге выбирает состояние с наибольшей оценкой: -7, затем -6, затем -5 и так далее, пока агент не достигнет цели.

Подведём итог:

\CourseFigure{vbm_1.jpg}{Обучение с подкреплением на основе ценности: обзор функции ценности}

\CourseFigure{vbm_2.jpg}{Методы на основе функции ценности: выбор действий по ценности состояний}

\section{Почему обучение с подкреплением называется «глубоким»?}\label{the-deep-in-reinforcement-learning}

\begin{quote}
\textbf{Примечание к курсу.} До сих пор мы говорили об обучении с подкреплением. Но откуда в названии взялось слово «глубокое»?
\end{quote}

Глубокое обучение с подкреплением предполагает \textbf{использование глубоких нейронных сетей для решения задач RL} --- отсюда и название.

Например, в следующем модуле мы изучим два алгоритма на основе функции ценности: Q-Learning (классическое RL), а затем Deep Q-Learning.

Вы увидите, что в первом случае \textbf{мы применяем традиционный алгоритм}, чтобы создать Q-таблицу, которая подсказывает действие для каждого состояния.

Во втором подходе \textbf{используется нейронная сеть} для аппроксимации Q-значений.

\CourseFigure{deep.jpg}{Схема по мотивам учебного ноутбука Udacity по Q-Learning}

Если вы ещё не знакомы с глубоким обучением, рекомендую посмотреть бесплатный курс \href{https://course.fast.ai}{FastAI Practical Deep Learning for Coders}.

\section{Краткое резюме}\label{summary}

Мы рассмотрели немало информации! Подведём итоги:

\begin{itemize}
\item
  Обучение с подкреплением --- вычислительный подход к обучению на основе действий. Мы создаём агента, который учится \textbf{взаимодействовать со средой методом проб и ошибок} и получает вознаграждения (отрицательные или положительные) в качестве обратной связи.
\item
  Цель любого RL-агента --- максимизировать ожидаемое суммарное вознаграждение (ожидаемую отдачу). В основе RL лежит \textbf{гипотеза вознаграждения}: \textbf{любую цель можно описать как максимизацию ожидаемого суммарного вознаграждения.}
\item
  Процесс RL --- цикл, порождающий последовательность \textbf{«состояние, действие, вознаграждение, следующее состояние».}
\item
  Чтобы рассчитать ожидаемую суммарную отдачу, мы дисконтируем вознаграждения: ближайшие по времени вознаграждения (в начале игры) \textbf{вероятнее получить, поскольку их проще предсказать, чем отдалённые будущие вознаграждения.}
\item
  Для решения задачи RL нужно \textbf{найти оптимальную политику}. Политика --- «мозг» агента, который определяет, \textbf{какое действие совершить в текущем состоянии}. Оптимальная политика \textbf{выбирает действия, максимизирующие ожидаемую отдачу.}
\item
  Оптимальную политику можно искать двумя способами:

  \begin{enumerate}
  \def\labelenumi{\arabic{enumi}.}
  \tightlist
  \item
    Непосредственно обучать политику: \textbf{методы на основе политики.}
  \item
    Обучать функцию ценности, оценивающую ожидаемую отдачу в каждом состоянии, и использовать её для выбора политики: \textbf{методы на основе функции ценности.}
  \end{enumerate}
\item
  Наконец, слово «глубокое» означает использование \textbf{глубоких нейронных сетей для оценки выбираемого действия (методы на основе политики) или ценности состояния (методы на основе функции ценности).}
\end{itemize}

\section{Глоссарий}\label{glossary}

Этот глоссарий создан сообществом. Мы приветствуем дополнения!

\subsubsection{Агент}\label{agent}

Агент учится \textbf{принимать решения методом проб и ошибок, получая от среды вознаграждения и наказания.}

\subsubsection{Среда}\label{environment}

Среда --- это моделируемый мир, \textbf{в котором агент может учиться, взаимодействуя с ним.}

\subsubsection{Марковское свойство}\label{markov-property-1}

Оно означает, что выбранное агентом действие \textbf{зависит только от текущего состояния и не зависит от предыдущих состояний и действий.}

\subsubsection{Наблюдения и состояния}\label{observationsstate}

\begin{itemize}
\tightlist
\item
  \textbf{Состояние}: полное описание состояния мира.
\item
  \textbf{Наблюдение}: частичное описание состояния среды или мира.
\end{itemize}

\subsubsection{Действия}\label{actions}

\begin{itemize}
\tightlist
\item
  \textbf{Дискретные действия}: конечное число действий, например влево, вправо, вверх и вниз.
\item
  \textbf{Непрерывные действия}: бесконечное число вариантов. Например, при управлении беспилотным автомобилем возможно бесконечно много управляющих воздействий.
\end{itemize}

\subsubsection{Вознаграждение и дисконтирование}\label{rewards-and-discounting}

\begin{itemize}
\tightlist
\item
  \textbf{Вознаграждение}: фундаментальный элемент RL; сообщает агенту, было ли действие хорошим или плохим.
\item
  Алгоритмы RL ориентированы на максимизацию \textbf{суммарного вознаграждения}.
\item
  \textbf{Гипотеза вознаграждения}: задачи RL можно сформулировать как максимизацию (суммарной) отдачи.
\item
  \textbf{Дисконтирование} применяется потому, что более ранние вознаграждения вероятнее получить: их проще предсказать, чем долгосрочные.
\end{itemize}

\subsubsection{Типы задач}\label{tasks}

\begin{itemize}
\tightlist
\item
  \textbf{Эпизодическая}: имеет начальную и конечную точки.
\item
  \textbf{Продолжающаяся}: имеет начало, но не имеет конечной точки.
\end{itemize}

\subsubsection{Компромисс между исследованием и использованием знаний}\label{exploration-vs-exploitation-trade-off}

\begin{itemize}
\tightlist
\item
  \textbf{Исследование}: выполнение случайных действий для изучения среды с получением обратной связи или вознаграждений.
\item
  \textbf{Использование знаний}: применение известных сведений о среде для получения максимального вознаграждения.
\item
  \textbf{Компромисс исследования и использования}: выбор баланса между \textbf{исследованием} среды и \textbf{использованием} уже имеющихся знаний.
\end{itemize}

\subsubsection{Политика}\label{policy}

\begin{itemize}
\tightlist
\item
  \textbf{Политика}: «мозг» агента, указывающий действие для заданного состояния.
\item
  \textbf{Оптимальная политика}: политика, \textbf{максимизирующая ожидаемую отдачу}, когда агент ей следует. Находится в процессе \emph{обучения}.
\end{itemize}

\subsubsection{Методы на основе политики}\label{policy-based-methods-1}

\begin{itemize}
\tightlist
\item
  Один из подходов к решению задач RL.
\item
  В нём политика обучается непосредственно.
\item
  Она сопоставляет каждому состоянию оптимальное действие либо распределение вероятностей по доступным в нём действиям.
\end{itemize}

\subsubsection{Методы на основе функции ценности}\label{value-based-methods-1}

\begin{itemize}
\tightlist
\item
  Ещё один подход к решению задач RL.
\item
  Здесь вместо обучения политики мы обучаем \textbf{функцию ценности}, сопоставляющую каждому состоянию ожидаемую ценность нахождения в нём.
\end{itemize}

Дополнения приветствуются!

Если вы хотите улучшить курс, можете \href{https://github.com/huggingface/deep-rl-class/pulls}{создать Pull Request}.

Этот глоссарий появился благодаря:

\begin{itemize}
\tightlist
\item
  \href{https://github.com/lucifermorningstar1305}{@lucifermorningstar1305}
\item
  \href{https://github.com/daspartho}{@daspartho}
\item
  \href{https://github.com/misza222}{@misza222}
\end{itemize}

\section{Обучаем первого агента глубокого обучения с подкреплением}\label{train-your-first-deep-reinforcement-learning-agent}

Исходный ноутбук: \url{https://colab.research.google.com/github/huggingface/deep-rl-class/blob/main/notebooks/unit1/unit1.ipynb}

Теперь, когда вы изучили основы обучения с подкреплением, можно обучить первого агента и поделиться результатом с сообществом через Hub: лунный посадочный модуль научится правильно садиться на поверхность Луны.

\CourseFigure{lunarLander-still.jpg}{Среда LunarLander: статичный кадр с официальной обложки курса; \href{https://huggingface.co/datasets/huggingface-deep-rl-course/course-images/resolve/main/en/unit1/lunarLander.gif}{оригинальная GIF-анимация}}

И наконец, вы \textbf{загрузите обученного агента в Hugging Face Hub --- бесплатную открытую платформу для обмена моделями, данными и демонстрациями.}

На нашей \href{https://huggingface.co/spaces/huggingface-projects/Deep-Reinforcement-Learning-Leaderboard}{таблице лидеров} вы сможете сравнить результаты с другими участниками и обменяться лучшими практиками повышения качества агента. Кто победит в соревновании первого модуля?

Чтобы это практическое задание было зачтено для \href{https://huggingface.co/deep-rl-course/en/unit0/introduction\#certification-process}{получения сертификата}, необходимо загрузить модель в Hub и \textbf{получить результат не ниже 200}.

Чтобы узнать свой результат, найдите модель в \href{https://huggingface.co/spaces/huggingface-projects/Deep-Reinforcement-Learning-Leaderboard}{таблице лидеров}; \textbf{результат = mean\_reward - std вознаграждения}.

\textbf{Если модели нет в списке, прокрутите страницу вниз и нажмите кнопку обновления.}

Подробнее о сертификации: \url{https://huggingface.co/deep-rl-course/en/unit0/introduction\#certification-process}

Отслеживать свой прогресс можно здесь: \url{https://huggingface.co/spaces/ThomasSimonini/Check-my-progress-Deep-RL-Course}

Приступим!

\textbf{Чтобы начать практическое задание, нажмите Open In Colab}:

\href{https://colab.research.google.com/github/huggingface/deep-rl-class/blob/master/notebooks/unit1/unit1.ipynb}{Открыть исходный ноутбук в Google Colab}

Мы настоятельно \textbf{рекомендуем студентам выполнять практические задания в Google Colab}, а не на личных компьютерах.

Благодаря Google Colab \textbf{вы можете сосредоточиться на обучении и экспериментах, не отвлекаясь на техническую настройку окружения.}

\section{Модуль 1: обучаем первого агента Deep RL}\label{unit-1-train-your-first-deep-reinforcement-learning-agent}

\CourseFigure{thumbnail.jpg}{Модуль 1: введение в глубокое обучение с подкреплением (официальная иллюстрация курса)}

В этом ноутбуке вы обучите своего \textbf{первого агента глубокого обучения с подкреплением} --- лунный посадочный модуль, который научится \textbf{правильно садиться на Луну}. Мы будем использовать библиотеку \href{https://stable-baselines3.readthedocs.io/en/master/}{Stable-Baselines3}, поделимся моделью с сообществом и поэкспериментируем с настройками.

\subsubsection{Среда}\label{the-environment}

\begin{itemize}
\tightlist
\item
  \href{https://gymnasium.farama.org/environments/box2d/lunar_lander/}{LunarLander-v2}
\end{itemize}

\subsubsection{Используемая библиотека}\label{the-library-used}

\begin{itemize}
\tightlist
\item
  \href{https://stable-baselines3.readthedocs.io/en/master/}{Stable-Baselines3}
\end{itemize}

Мы постоянно стараемся улучшать наши уроки. Если \textbf{вы обнаружили проблему в этом ноутбуке}, пожалуйста, \href{https://github.com/huggingface/deep-rl-class/issues}{создайте Issue в репозитории GitHub}.

\subsection{Цели ноутбука}\label{objectives-of-this-notebook}

После выполнения ноутбука вы сможете:

\begin{itemize}
\tightlist
\item
  Пользоваться \textbf{Gymnasium} --- библиотекой сред.
\item
  Пользоваться \textbf{Stable-Baselines3} --- библиотекой глубокого обучения с подкреплением.
\item
  \textbf{Публиковать обученного агента в Hub} вместе с видеозаписью его работы и оценкой качества.
\end{itemize}

\subsection{Ноутбук курса Deep Reinforcement Learning}\label{this-notebook-is-from-deep-reinforcement-learning-course}

\CourseFigure{deep-rl-course-illustration.jpg}{Обзор курса Deep Reinforcement Learning}

В этом бесплатном курсе вы:

\begin{itemize}
\tightlist
\item
  Изучите Deep RL \textbf{в теории и на практике}.
\item
  Научитесь \textbf{использовать популярные библиотеки Deep RL}: Stable Baselines3, RL Baselines3 Zoo, CleanRL и Sample Factory 2.0.
\item
  Обучите \textbf{агентов в необычных средах}.
\item
  \textbf{Получите сертификат об окончании курса}, выполнив 80\% заданий.
\end{itemize}

И это ещё не всё!

Программа курса: \url{https://simoninithomas.github.io/deep-rl-course}

Не забудьте \textbf{\href{http://eepurl.com/ic5ZUD}{зарегистрироваться на курс}} (мы собираем адреса электронной почты, чтобы \textbf{присылать ссылки по мере публикации модулей, а также информацию о заданиях и обновлениях}).

Лучший способ оставаться на связи и задавать вопросы --- \textbf{присоединиться к нашему Discord-серверу}, чтобы общаться с сообществом и преподавателями: \url{https://discord.gg/ydHrjt3WP5}

\subsection{Предварительные требования}\label{prerequisites}

Перед выполнением ноутбука необходимо:

\textbf{\href{https://huggingface.co/deep-rl-course/unit0/introduction}{Прочитать модуль 0}}, где собрана \textbf{основная информация о курсе и его организации}.

\textbf{Разобраться в основах обучения с подкреплением} (MC, TD, гипотеза вознаграждения и др.), \href{https://huggingface.co/deep-rl-course/unit1/introduction}{прочитав модуль 1}.

\subsection{Краткое повторение Deep RL}\label{a-small-recap-of-deep-reinforcement-learning}

\CourseFigure{RL_process_game.jpg}{Цикл взаимодействия в RL на примере видеоигры}

Кратко повторим изученное в первом модуле:

\begin{itemize}
\item
  Обучение с подкреплением --- \textbf{вычислительный подход к обучению на основе действий}. Мы создаём агента, который \textbf{взаимодействует со средой методом проб и ошибок} и получает вознаграждения (отрицательные или положительные) в качестве обратной связи.
\item
  Цель любого RL-агента --- \textbf{максимизировать ожидаемое суммарное вознаграждение} (ожидаемую отдачу), поскольку RL основано на \emph{гипотезе вознаграждения}: любую цель можно описать как максимизацию ожидаемого суммарного вознаграждения.
\item
  Процесс RL --- \textbf{цикл, создающий последовательность состояний, действий, вознаграждений и следующих состояний}.
\item
  Для расчёта ожидаемого суммарного вознаграждения \textbf{мы дисконтируем награды}: более близкие по времени вознаграждения вероятнее получить, поскольку их проще предсказать, чем далёкие будущие.
\item
  Чтобы решить задачу RL, необходимо \textbf{найти оптимальную политику}; политика --- «мозг» агента, указывающий действие для текущего состояния. Оптимальная политика максимизирует ожидаемую отдачу.
\end{itemize}

Оптимальную политику можно найти \textbf{двумя} способами:

\begin{itemize}
\item
  \textbf{Непосредственно обучать политику}: методы на основе политики.
\item
  \textbf{Обучать функцию ценности}, оценивающую ожидаемую отдачу в каждом состоянии, и использовать её для построения политики: методы на основе функции ценности.
\item
  Наконец, глубоким RL называют потому, что \textbf{для оценки выбираемого действия (методы на основе политики) или ценности состояния (методы на основе функции ценности) мы используем глубокие нейронные сети.}
\end{itemize}

\section{Обучим первого агента Deep RL и загрузим его в Hub}\label{lets-train-our-first-deep-reinforcement-learning-agent-and-upload-it-to-the-hub}

\subsection{Получение сертификата}\label{get-a-certificate}

Чтобы практическое задание было засчитано при \href{https://huggingface.co/deep-rl-course/en/unit0/introduction\#certification-process}{сертификации}, загрузите обученную модель в Hub и \textbf{получите результат не ниже 200}.

Чтобы узнать результат, найдите модель в \href{https://huggingface.co/spaces/huggingface-projects/Deep-Reinforcement-Learning-Leaderboard}{таблице лидеров}: \textbf{результат = mean\_reward - стандартное отклонение вознаграждения}.

Дополнительная информация о сертификации: \url{https://huggingface.co/deep-rl-course/en/unit0/introduction\#certification-process}

\subsection{Настройка GPU}\label{set-the-gpu}

\begin{itemize}
\tightlist
\item
  Для \textbf{ускорения обучения агента воспользуемся GPU}. Для этого в Colab выберите \passthrough{\lstinline!Runtime > Change Runtime type!}.
\end{itemize}

\CourseFigure{gpu-step1.jpg}{Меню Google Colab для изменения типа среды выполнения}

\begin{itemize}
\tightlist
\item
  Выберите \passthrough{\lstinline!Hardware Accelerator > GPU!}.
\end{itemize}

\CourseFigure{gpu-step2.jpg}{Выбор аппаратного ускорителя GPU в Google Colab}

\subsection{Установка зависимостей и создание виртуального экрана}\label{install-dependencies-and-create-a-virtual-screen}

Сначала установим необходимые зависимости: нам потребуется несколько библиотек.

\begin{itemize}
\tightlist
\item
  \passthrough{\lstinline!gymnasium[box2d]!}: содержит среду LunarLander-v2.
\item
  \passthrough{\lstinline!stable-baselines3[extra]!}: библиотека глубокого обучения с подкреплением.
\item
  \passthrough{\lstinline!huggingface_sb3!}: дополнительные средства Stable-Baselines3 для загрузки моделей в Hugging Face Hub и скачивания их оттуда.
\end{itemize}

Чтобы упростить установку, мы создали скрипт, устанавливающий все эти зависимости.

\begin{lstlisting}[language=bash]
apt install swig cmake
\end{lstlisting}

\begin{lstlisting}[language=bash]
pip install -r https://raw.githubusercontent.com/huggingface/deep-rl-class/main/notebooks/unit1/requirements-unit1.txt
\end{lstlisting}

Во время работы с ноутбуком потребуется записать видео поведения агента. Для этого в Colab \textbf{нужно создать виртуальный экран, позволяющий визуализировать среду} и записывать кадры.

Поэтому следующая ячейка устанавливает библиотеки виртуального экрана, создаёт его и запускает.

\begin{lstlisting}[language=bash]
sudo apt-get update
apt install python3-opengl
apt install ffmpeg
apt install xvfb
pip3 install pyvirtualdisplay
\end{lstlisting}

Чтобы новые библиотеки начали использоваться, \textbf{иногда необходимо перезапустить среду выполнения ноутбука}. Следующая ячейка намеренно вызовет \textbf{аварийную остановку среды, после чего понадобится снова подключиться и выполнять код, начиная с этого места}. Такой приём позволит запустить виртуальный экран.

\begin{lstlisting}[language=Python]
import os

os.kill(os.getpid(), 9)
\end{lstlisting}

\begin{lstlisting}[language=Python]
# Virtual display
from pyvirtualdisplay import Display

virtual_display = Display(visible=0, size=(1400, 900))
virtual_display.start()
\end{lstlisting}

\subsection{Импорт библиотек}\label{import-the-packages}

Дополнительно импортируем библиотеку huggingface\_hub, \textbf{чтобы загружать обученные модели в Hub и скачивать их оттуда}.

Hugging Face Hub --- единая платформа, где любой желающий может публиковать и изучать модели и наборы данных. Здесь есть версионирование, метрики, визуализации и другие функции, упрощающие совместную работу.

Посмотреть доступные модели глубокого обучения с подкреплением можно по ссылке: \url{https://huggingface.co/models?pipeline_tag=reinforcement-learning&sort=downloads}

\begin{lstlisting}[language=Python]
import gymnasium

from huggingface_sb3 import load_from_hub, package_to_hub
from huggingface_hub import (
    notebook_login,
)  # To log to our Hugging Face account to be able to upload models to the Hub.

from stable_baselines3 import PPO
from stable_baselines3.common.env_util import make_vec_env
from stable_baselines3.common.evaluation import evaluate_policy
from stable_baselines3.common.monitor import Monitor
\end{lstlisting}

\subsection{Как работает Gymnasium}\label{understand-gymnasium-and-how-it-works}

Библиотека, содержащая нашу среду, называется Gymnasium. \textbf{При изучении глубокого RL вы будете часто пользоваться Gymnasium.}

Gymnasium --- \textbf{новая версия библиотеки Gym}, \href{https://farama.org/}{поддерживаемая организацией Farama Foundation}.

Gymnasium предоставляет две основные возможности:

\begin{itemize}
\tightlist
\item
  Интерфейс, позволяющий \textbf{создавать среды RL}.
\item
  \textbf{Набор готовых сред} (gym-control, Atari, Box2D и т. д.).
\end{itemize}

Рассмотрим пример, но сначала вспомним цикл RL.

\CourseFigure{RL_process_game.jpg}{Цикл взаимодействия в RL на примере видеоигры}

На каждом шаге:

\begin{itemize}
\tightlist
\item
  Агент получает от \textbf{среды} \textbf{состояние (S0)} --- первый кадр игры.
\item
  На основе \textbf{состояния (S0)} агент выбирает \textbf{действие (A0)} --- движение вправо.
\item
  Среда переходит в \textbf{новое состояние (S1)} --- следующий кадр.
\item
  Среда выдаёт агенту \textbf{вознаграждение (R1)} --- мы не погибли, то есть получили \emph{положительное вознаграждение +1}.
\end{itemize}

В Gymnasium это выглядит так:

1 Создаём среду с помощью \passthrough{\lstinline!gymnasium.make()!}.

2 Сбрасываем среду в начальное состояние командой \passthrough{\lstinline!observation = env.reset()!}.

На каждом шаге:

3 Получаем действие от модели (в нашем примере выбирается случайное действие).

4 С помощью \passthrough{\lstinline!env.step(action)!} выполняем действие в среде и получаем:

\begin{itemize}
\tightlist
\item
  \passthrough{\lstinline!observation!}: новое состояние (st+1).
\item
  \passthrough{\lstinline!reward!}: вознаграждение после выполнения действия.
\item
  \passthrough{\lstinline!terminated!}: признак завершения эпизода (агент достиг терминального состояния).
\item
  \passthrough{\lstinline!truncated!}: новый для этой версии признак обрыва эпизода, например по ограничению времени или при выходе агента за границы среды.
\item
  \passthrough{\lstinline!info!}: словарь с дополнительной информацией (зависит от среды).
\end{itemize}

Подробнее см. \url{https://gymnasium.farama.org/api/env/\#gymnasium.Env.step}

Если эпизод завершился:

\begin{itemize}
\tightlist
\item
  Сбрасываем среду в начальное состояние командой \passthrough{\lstinline!observation = env.reset()!}.
\end{itemize}

\textbf{Рассмотрим пример!} Обязательно прочитайте код.

\begin{lstlisting}[language=Python]
import gymnasium as gym

# First, we create our environment called LunarLander-v2
env = gym.make("LunarLander-v2")

# Then we reset this environment
observation, info = env.reset()

for _ in range(20):
    # Take a random action
    action = env.action_space.sample()
    print("Action taken:", action)

    # Do this action in the environment and get
    # next_state, reward, terminated, truncated and info
    observation, reward, terminated, truncated, info = env.step(action)

    # If the game is terminated (in our case we land, crashed) or truncated (timeout)
    if terminated or truncated:
        # Reset the environment
        print("Environment is reset")
        observation, info = env.reset()

env.close()
\end{lstlisting}

\subsection{Создание среды LunarLander и знакомство с ней}\label{create-the-lunarlander-environment-and-understand-how-it-works}

\subsubsection{Среда}\label{the-environment-1}

В этом первом руководстве мы научим агента \href{https://gymnasium.farama.org/environments/box2d/lunar_lander/}{Lunar Lander} \textbf{правильно садиться на Луну}. Для этого он должен научиться \textbf{управлять скоростью и положением (по горизонтали, вертикали и углу наклона), чтобы выполнить безопасную посадку.}

\begin{center}\rule{0.5\linewidth}{0.5pt}\end{center}

При работе с новой средой полезно сначала изучить её документацию:

\url{https://gymnasium.farama.org/environments/box2d/lunar_lander/}

\begin{center}\rule{0.5\linewidth}{0.5pt}\end{center}

Посмотрим, как выглядит эта среда:

\begin{lstlisting}[language=Python]
# We create our environment with gym.make("<name_of_the_environment>")
env = gym.make("LunarLander-v2")
env.reset()
print("_____OBSERVATION SPACE_____ \n")
print("Observation Space Shape", env.observation_space.shape)
print("Sample observation", env.observation_space.sample())  # Get a random observation
\end{lstlisting}

Из строки \passthrough{\lstinline!Observation Space Shape (8,)!} видно, что наблюдение --- вектор из 8 элементов, каждый из которых содержит определённую информацию о посадочном модуле:

\begin{itemize}
\tightlist
\item
  Горизонтальная координата (x).
\item
  Вертикальная координата (y).
\item
  Горизонтальная скорость (x).
\item
  Вертикальная скорость (y).
\item
  Угол наклона.
\item
  Угловая скорость.
\item
  Признак касания земли левой опорой (логический).
\item
  Признак касания земли правой опорой (логический).
\end{itemize}

\begin{lstlisting}[language=Python]
print("\n _____ACTION SPACE_____ \n")
print("Action Space Shape", env.action_space.n)
print("Action Space Sample", env.action_space.sample())  # Take a random action
\end{lstlisting}

Пространство действий (множество доступных агенту действий) дискретно и состоит из четырёх вариантов:

\begin{itemize}
\tightlist
\item
  Действие 0: ничего не делать.
\item
  Действие 1: включить левый боковой двигатель ориентации.
\item
  Действие 2: включить главный двигатель.
\item
  Действие 3: включить правый боковой двигатель ориентации.
\end{itemize}

Функция вознаграждения (выдающая награду на каждом временном шаге):

После каждого шага агент получает вознаграждение. Суммарное вознаграждение за эпизод --- \textbf{сумма вознаграждений за все шаги этого эпизода}.

На каждом шаге вознаграждение:

\begin{itemize}
\tightlist
\item
  Увеличивается или уменьшается по мере приближения или удаления от посадочной площадки.
\item
  Увеличивается при замедлении и уменьшается при ускорении посадочного модуля.
\item
  Уменьшается при увеличении наклона модуля (отклонения от горизонтали).
\item
  Увеличивается на 10 баллов за каждую опору, касающуюся земли.
\item
  Уменьшается на 0,03 балла за каждый кадр работы бокового двигателя.
\item
  Уменьшается на 0,3 балла за каждый кадр работы главного двигателя.
\end{itemize}

Кроме того, эпизод приносит \textbf{дополнительное вознаграждение -100 при аварии или +100 при успешной посадке}.

Эпизод \textbf{считается успешным, если итоговый результат составляет не менее 200 баллов}.

\paragraph{Векторизованная среда}\label{vectorized-environment}

\begin{itemize}
\tightlist
\item
  Создадим векторизованную среду (способ объединить несколько независимых сред в одну) из 16 экземпляров. Так \textbf{во время обучения агент получит более разнообразный опыт}.
\end{itemize}

\begin{lstlisting}[language=Python]
# Create the environment
env = make_vec_env("LunarLander-v2", n_envs=16)
\end{lstlisting}

\subsection{Создание модели}\label{create-the-model}

\begin{itemize}
\item
  Мы изучили среду и поняли задачу: \textbf{научиться безопасно сажать LunarLander на посадочную площадку, управляя левым и правым боковыми и главным двигателями}. Теперь построим алгоритм решения.
\item
  Для этого используем первую библиотеку Deep RL в курсе --- \href{https://stable-baselines3.readthedocs.io/en/master/}{Stable Baselines3 (SB3)}.
\item
  SB3 --- набор \textbf{надёжных реализаций алгоритмов обучения с подкреплением на PyTorch}.
\end{itemize}

\begin{center}\rule{0.5\linewidth}{0.5pt}\end{center}

При знакомстве с новой библиотекой полезно сначала изучить \href{https://stable-baselines3.readthedocs.io/en/master/}{документацию}, а затем выполнить несколько руководств.

\begin{center}\rule{0.5\linewidth}{0.5pt}\end{center}

\CourseFigure{sb3.png}{Иллюстрация библиотеки Stable-Baselines3}

Для решения задачи воспользуемся алгоритмом \textbf{PPO} из SB3. \href{https://stable-baselines3.readthedocs.io/en/master/modules/ppo.html\#example\%5D}{PPO (Proximal Policy Optimization) --- один из современных алгоритмов глубокого обучения с подкреплением, которые будут рассмотрены в этом курсе}.

В PPO объединяются:

\begin{itemize}
\tightlist
\item
  \emph{Метод RL на основе функции ценности}: обучение функции ценности действия, показывающей \textbf{наиболее ценное действие для заданной пары состояния и действия}.
\item
  \emph{Метод RL на основе политики}: обучение политики, \textbf{возвращающей распределение вероятностей по действиям}.
\end{itemize}

Настроить Stable-Baselines3 достаточно просто:

1 \textbf{Создайте среду} (в нашем случае это уже сделано выше).

2 \textbf{Выберите модель и создайте её экземпляр}: \passthrough{\lstinline!model = PPO("MlpPolicy")!}.

3 \textbf{Обучите агента} через \passthrough{\lstinline!model.learn!}, задав число временных шагов обучения.

\begin{lstlisting}
# Create environment
env = gym.make('LunarLander-v2')

# Instantiate the agent
model = PPO('MlpPolicy', env, verbose=1)
# Train the agent
model.learn(total_timesteps=int(2e5))
\end{lstlisting}

\begin{lstlisting}[language=Python]
# TODO: Define a PPO MlpPolicy architecture
# We use MultiLayerPerceptron (MLPPolicy) because the input is a vector,
# if we had frames as input we would use CnnPolicy
model =
\end{lstlisting}

\paragraph{Решение}\label{solution}

\begin{lstlisting}[language=Python]
# SOLUTION
# We added some parameters to accelerate the training
model = PPO(
    policy="MlpPolicy",
    env=env,
    n_steps=1024,
    batch_size=64,
    n_epochs=4,
    gamma=0.999,
    gae_lambda=0.98,
    ent_coef=0.01,
    verbose=1,
)
\end{lstlisting}

\subsection{Обучение агента PPO}\label{train-the-ppo-agent}

\begin{itemize}
\tightlist
\item
  Обучим агента в течение 1 000 000 шагов; не забудьте включить GPU в Colab. Это займёт примерно \textasciitilde20 минут, но для первого эксперимента можно использовать меньше шагов.
\item
  Пока идёт обучение, сделайте заслуженный перерыв!
\end{itemize}

\begin{lstlisting}[language=Python]
# TODO: Train it for 1,000,000 timesteps
# TODO: Specify file name for model and save the model to file
model_name = "ppo-LunarLander-v2"
\end{lstlisting}

\paragraph{Решение}\label{solution-1}

\begin{lstlisting}[language=Python]
# SOLUTION
# Train it for 1,000,000 timesteps
model.learn(total_timesteps=1000000)
# Save the model
model_name = "ppo-LunarLander-v2"
model.save(model_name)
\end{lstlisting}

\subsection{Оценка качества агента}\label{evaluate-the-agent}

\begin{itemize}
\tightlist
\item
  Не забудьте обернуть среду в \href{https://stable-baselines3.readthedocs.io/en/master/common/monitor.html}{Monitor}.
\item
  Теперь, когда агент LunarLander обучен, нужно \textbf{проверить качество его работы}.
\item
  Для этого Stable-Baselines3 предоставляет функцию \passthrough{\lstinline!evaluate_policy!}.
\item
  Чтобы заполнить этот раздел, \href{https://stable-baselines3.readthedocs.io/en/master/guide/examples.html\#basic-usage-training-saving-loading}{обратитесь к документации}.
\item
  Далее мы узнаем, \textbf{как автоматически оценивать и публиковать агента для соревнования в таблице лидеров, но сначала сделаем оценку вручную}.
\end{itemize}

При оценке качества следует создавать отдельную среду для тестирования, а не использовать обучающую.

\begin{lstlisting}[language=Python]
# TODO: Evaluate the agent
# Create a new environment for evaluation
eval_env =

# Evaluate the model with 10 evaluation episodes and deterministic=True
mean_reward, std_reward = 

# Print the results
\end{lstlisting}

\paragraph{Решение}\label{solution-2}

\begin{lstlisting}[language=Python]
# @title
eval_env = Monitor(gym.make("LunarLander-v2"))
mean_reward, std_reward = evaluate_policy(model, eval_env, n_eval_episodes=10, deterministic=True)
print(f"mean_reward={mean_reward:.2f} +/- {std_reward}")
\end{lstlisting}

\begin{itemize}
\tightlist
\item
  В моём случае после миллиона шагов среднее вознаграждение составило \passthrough{\lstinline!200.20 +/- 20.80!}. Это означает, что наш лунный модуль готов к посадке на Луну.
\end{itemize}

\subsection{Публикация обученной модели в Hub}\label{publish-our-trained-model-on-the-hub}

Теперь, когда мы получили хорошие результаты, можно опубликовать обученную модель в Hub одной строкой кода.

Документация библиотеки: \url{https://github.com/huggingface/huggingface_sb3/tree/main\#hugging-face--x-stable-baselines3-v20}

Вот пример карточки модели (для Space Invaders):

С помощью \passthrough{\lstinline!package_to_hub!} \textbf{можно оценить агента, записать видео его работы, создать карточку модели и загрузить всё в Hub}.

Таким образом:

\begin{itemize}
\tightlist
\item
  Вы можете \textbf{демонстрировать свою работу}.
\item
  Вы можете \textbf{наблюдать за тем, как играет агент}.
\item
  Вы можете \textbf{поделиться с сообществом агентом, которого смогут использовать другие}.
\item
  Вы можете \textbf{сравнивать качество своего агента с результатами других участников в таблице лидеров}: \url{https://huggingface.co/spaces/huggingface-projects/Deep-Reinforcement-Learning-Leaderboard}
\end{itemize}

Чтобы опубликовать модель для сообщества, нужно выполнить ещё три шага:

1 Если ещё не сделали этого, создайте аккаунт Hugging Face: \url{https://huggingface.co/join}

2 Войдите в аккаунт и сохраните токен аутентификации с сайта Hugging Face.

\begin{itemize}
\tightlist
\item
  Создайте новый токен (\url{https://huggingface.co/settings/tokens}) \textbf{с правом записи (write)}.
\end{itemize}

\CourseFigure{create-token.jpg}{Создание токена Hugging Face Hub с правами на запись}

\begin{itemize}
\tightlist
\item
  Скопируйте токен.
\item
  Выполните ячейку ниже и вставьте токен.
\end{itemize}

\begin{lstlisting}[language=Python]
notebook_login()
!git config --global credential.helper store
\end{lstlisting}

Если вы не используете Google Colab или Jupyter Notebook, вместо этого выполните команду \passthrough{\lstinline!huggingface-cli login!}.

3 Теперь всё готово к публикации обученного агента в Hub через функцию \passthrough{\lstinline!package_to_hub()!}.

Заполним параметры функции \passthrough{\lstinline!package_to_hub!}:

\begin{itemize}
\tightlist
\item
  \passthrough{\lstinline!model!}: наша обученная модель.
\item
  \passthrough{\lstinline!model_name!}: имя обученной модели, которое мы задали в \passthrough{\lstinline!model_save!}.
\item
  \passthrough{\lstinline!model_architecture!}: архитектура модели, в нашем случае PPO.
\item
  \passthrough{\lstinline!env_id!}: название среды, в нашем случае \passthrough{\lstinline!LunarLander-v2!}.
\item
  \passthrough{\lstinline!eval_env!}: среда оценки, заданная переменной eval\_env.
\item
  \passthrough{\lstinline!repo_id!}: имя репозитория Hugging Face Hub, который будет создан или обновлён, \passthrough{\lstinline!(repo_id = \{username\}/\{repo_name\})!}.
\end{itemize}

\textbf{Подходящий формат имени --- \passthrough{\lstinline!\{username\}/\{model_architecture\}-\{env_id\}!}}.

\begin{itemize}
\tightlist
\item
  \passthrough{\lstinline!commit_message!}: сообщение коммита.
\end{itemize}

\begin{lstlisting}[language=Python]
import gymnasium as gym
from stable_baselines3.common.vec_env import DummyVecEnv
from stable_baselines3.common.env_util import make_vec_env

from huggingface_sb3 import package_to_hub

## TODO: Define a repo_id
## repo_id is the id of the model repository from the Hugging Face Hub (repo_id = {organization}/{repo_name} for instance ThomasSimonini/ppo-LunarLander-v2
repo_id = 

# TODO: Define the name of the environment
env_id = 

# Create the evaluation env and set the render_mode="rgb_array"
eval_env = DummyVecEnv([lambda: gym.make(env_id, render_mode="rgb_array")])


# TODO: Define the model architecture we used
model_architecture = ""

## TODO: Define the commit message
commit_message = ""

# method save, evaluate, generate a model card and record a replay video of your agent before pushing the repo to the hub
package_to_hub(model=model, # Our trained model
               model_name=model_name, # The name of our trained model 
               model_architecture=model_architecture, # The model architecture we used: in our case PPO
               env_id=env_id, # Name of the environment
               eval_env=eval_env, # Evaluation Environment
               repo_id=repo_id, # id of the model repository from the Hugging Face Hub (repo_id = {organization}/{repo_name} for instance ThomasSimonini/ppo-LunarLander-v2
               commit_message=commit_message)
\end{lstlisting}

\paragraph{Решение}\label{solution-3}

\begin{lstlisting}[language=Python]
import gymnasium as gym

from stable_baselines3 import PPO
from stable_baselines3.common.vec_env import DummyVecEnv
from stable_baselines3.common.env_util import make_vec_env

from huggingface_sb3 import package_to_hub

# PLACE the variables you've just defined two cells above
# Define the name of the environment
env_id = "LunarLander-v2"

# TODO: Define the model architecture we used
model_architecture = "PPO"

## Define a repo_id
## repo_id is the id of the model repository from the Hugging Face Hub (repo_id = {organization}/{repo_name} for instance ThomasSimonini/ppo-LunarLander-v2
## CHANGE WITH YOUR REPO ID
repo_id = "ThomasSimonini/ppo-LunarLander-v2"  # Change with your repo id, you can't push with mine :)

## Define the commit message
commit_message = "Upload PPO LunarLander-v2 trained agent"

# Create the evaluation env and set the render_mode="rgb_array"
eval_env = DummyVecEnv([lambda: Monitor(gym.make(env_id, render_mode="rgb_array"))])

# PLACE the package_to_hub function you've just filled here
package_to_hub(
    model=model,  # Our trained model
    model_name=model_name,  # The name of our trained model
    model_architecture=model_architecture,  # The model architecture we used: in our case PPO
    env_id=env_id,  # Name of the environment
    eval_env=eval_env,  # Evaluation Environment
    repo_id=repo_id,  # id of the model repository from the Hugging Face Hub (repo_id = {organization}/{repo_name} for instance ThomasSimonini/ppo-LunarLander-v2
    commit_message=commit_message,
)
\end{lstlisting}

Поздравляем! Вы обучили и опубликовали своего первого агента глубокого RL. Скрипт выше должен вывести ссылку на репозиторий модели, например \url{https://huggingface.co/osanseviero/test_sb3}. На этой странице можно:

\begin{itemize}
\tightlist
\item
  Посмотреть справа видеопревью работы агента.
\item
  Перейти в раздел «Files and versions», чтобы увидеть все файлы репозитория.
\item
  Нажать «Use in stable-baselines3», чтобы получить фрагмент кода для загрузки модели.
\item
  Прочитать карточку модели (файл \passthrough{\lstinline!README.md!}), содержащую описание модели.
\end{itemize}

В основе Hub лежат репозитории Git (не страшно, если вы пока не знакомы с Git). Это позволяет загружать новые версии модели по мере экспериментов и улучшения качества агента.

Сравните результаты LunarLander-v2 с результатами других студентов с помощью \href{https://huggingface.co/spaces/huggingface-projects/Deep-Reinforcement-Learning-Leaderboard}{таблицы лидеров}.

\subsection{Загрузка сохранённой модели LunarLander из Hub}\label{load-a-saved-lunarlander-model-from-the-hub}

Благодарим \href{https://github.com/ironbar}{ironbar} за вклад в материалы.

Загрузить сохранённую модель из Hub очень просто.

На странице \url{https://huggingface.co/models?library=stable-baselines3} можно увидеть все сохранённые модели Stable-Baselines3.

\begin{enumerate}
\def\labelenumi{\arabic{enumi}.}
\tightlist
\item
  Выберите одну из них и скопируйте её repo\_id.
\end{enumerate}

\CourseFigure{copy-id.png}{Копирование идентификатора репозитория модели в Hugging Face Hub}

\begin{enumerate}
\def\labelenumi{\arabic{enumi}.}
\setcounter{enumi}{1}
\tightlist
\item
  Затем используйте load\_from\_hub с параметрами:
\end{enumerate}

\begin{itemize}
\tightlist
\item
  Идентификатор repo\_id.
\item
  Имя файла: сохранённая модель в репозитории с расширением (*.zip).
\end{itemize}

Поскольку модель, которую я загружаю из Hub, была обучена с Gym (старой версией Gymnasium), нужно установить Shimmy --- средство преобразования API, позволяющее корректно запустить среду.

Документация Shimmy: \url{https://github.com/Farama-Foundation/Shimmy}

\begin{lstlisting}[language=Python]
!pip install shimmy
\end{lstlisting}

\begin{lstlisting}[language=Python]
from huggingface_sb3 import load_from_hub

repo_id = "Classroom-workshop/assignment2-omar"  # The repo_id
filename = "ppo-LunarLander-v2.zip"  # The model filename.zip

# When the model was trained on Python 3.8 the pickle protocol is 5
# But Python 3.6, 3.7 use protocol 4
# In order to get compatibility we need to:
# 1. Install pickle5 (we done it at the beginning of the colab)
# 2. Create a custom empty object we pass as parameter to PPO.load()
custom_objects = {
    "learning_rate": 0.0,
    "lr_schedule": lambda _: 0.0,
    "clip_range": lambda _: 0.0,
}

checkpoint = load_from_hub(repo_id, filename)
model = PPO.load(checkpoint, custom_objects=custom_objects, print_system_info=True)
\end{lstlisting}

Проверим качество этого агента:

\begin{lstlisting}[language=Python]
# @title
eval_env = Monitor(gym.make("LunarLander-v2"))
mean_reward, std_reward = evaluate_policy(model, eval_env, n_eval_episodes=10, deterministic=True)
print(f"mean_reward={mean_reward:.2f} +/- {std_reward}")
\end{lstlisting}

\subsection{Дополнительные задания}\label{some-additional-challenges}

Лучший способ учиться --- \textbf{экспериментировать самостоятельно}! Как вы видели, текущий агент работает не слишком хорошо. Сначала попробуйте увеличить число шагов обучения. При 1 000 000 шагов нам удалось получить отличные результаты!

Своих агентов можно найти в \href{https://huggingface.co/spaces/huggingface-projects/Deep-Reinforcement-Learning-Leaderboard}{таблице лидеров}. Сможете подняться на вершину рейтинга?

Вот несколько идей, как добиться этого:

\begin{itemize}
\tightlist
\item
  Увеличьте число шагов обучения.
\item
  Попробуйте другие гиперпараметры \passthrough{\lstinline!PPO!}; список приведён в \url{https://stable-baselines3.readthedocs.io/en/master/modules/ppo.html\#parameters}.
\item
  Ознакомьтесь с \href{https://stable-baselines3.readthedocs.io/en/master/modules/dqn.html}{документацией Stable-Baselines3} и попробуйте другую модель, например DQN.
\item
  \textbf{Загрузите новую обученную модель} в Hub.
\end{itemize}

\textbf{Сравните результаты своего LunarLander-v2 с результатами других студентов} через \href{https://huggingface.co/spaces/huggingface-projects/Deep-Reinforcement-Learning-Leaderboard}{таблицу лидеров}.

Посадка на Луну кажется слишком скучной? Попробуйте \textbf{сменить среду}: например, MountainCar-v0, CartPole-v1 или CarRacing-v0. Изучите их в \href{https://www.gymlibrary.dev/}{документации Gym} и экспериментируйте!

\begin{center}\rule{0.5\linewidth}{0.5pt}\end{center}

Поздравляем с завершением главы! Это была самая большая глава курса, \textbf{и информации действительно было много.}

Если какие-то понятия всё ещё вызывают затруднения, это совершенно нормально! \textbf{У меня и у всех изучавших RL было то же самое.}

Прежде чем двигаться дальше, постарайтесь \textbf{хорошо усвоить материал и выполнить дополнительные задания}. Важно освоить эти основы и приобрести прочную базу.

Конечно, далее мы будем изучать эти понятия подробнее, но \textbf{лучше разобраться с ними уже сейчас, прежде чем переходить к следующим главам}.

В следующем, первом бонусном модуле, вы научите пса Huggy приносить палку.

\CourseFigure{huggy-notebook.jpg}{Huggy в виртуальной обучающей среде}

\subsection{Продолжайте учиться, у вас всё получится}\label{keep-learning-stay-awesome}

\section{Тест для самопроверки}\label{quiz}

Лучший способ учиться и \href{https://www.coursera.org/lecture/learning-how-to-learn/illusions-of-competence-BuFzf}{не поддаваться иллюзии компетентности} --- \textbf{проверять себя}. Это поможет понять, \textbf{какие знания нужно закрепить}.

\subsubsection{Вопрос 1: что такое обучение с подкреплением?}\label{q1-what-is-reinforcement-learning}

\subsubsection{Вопрос 2: опишите цикл RL}\label{q2-define-the-rl-loop}

\CourseFigure{rl-loop-ex.jpg}{Задание теста: впишите состояние, действие, следующее состояние и вознаграждение в цикл RL}

На каждом шаге:

\begin{itemize}
\tightlist
\item
  Агент получает от среды \_\_\_\_\_\_.
\item
  На основе \_\_\_\_\_\_ агент выбирает \_\_\_\_\_\_.
\item
  Агент движется вправо.
\item
  Среда переходит в \_\_\_\_\_\_.
\item
  Среда выдаёт агенту \_\_\_\_\_\_.
\end{itemize}

\textbf{Выберите все правильные варианты (их может быть несколько).}

\begin{itemize}
\tightlist
\item
  \textbf{A.} действие a0, действие a0, состояние s0, состояние s1, вознаграждение r1.
\item
  \textbf{B.} состояние s0, состояние s0, действие a0, новое состояние s1, вознаграждение r1.
\item
  \textbf{C.} состояние s0, состояние s0, действие a0, состояние s1, действие a1.
\end{itemize}

\subsubsection{Вопрос 3: чем состояние отличается от наблюдения?}\label{q3-whats-the-difference-between-a-state-and-an-observation}

\textbf{Выберите все правильные варианты (их может быть несколько).}

\begin{itemize}
\tightlist
\item
  \textbf{A.} Состояние полностью описывает мир (скрытой информации нет).
\item
  \textbf{B.} Состояние представляет собой частичное описание мира.
\item
  \textbf{C.} Наблюдение полностью описывает мир (скрытой информации нет).
\item
  \textbf{D.} Наблюдение представляет собой частичное описание состояния.
\item
  \textbf{E.} В шахматной среде мы получаем состояние.
\item
  \textbf{F.} В шахматной среде мы получаем наблюдение.
\item
  \textbf{G.} В Super Mario Bros мы получаем состояние.
\item
  \textbf{H.} В Super Mario Bros мы получаем наблюдение.
\end{itemize}

\subsubsection{Вопрос 4: какие существуют два типа задач RL?}\label{q4-a-task-is-an-instance-of-a-reinforcement-learning-problem.-what-are-the-two-types-of-tasks}

\textbf{Выберите все правильные варианты (их может быть несколько).}

\begin{itemize}
\tightlist
\item
  \textbf{A.} Эпизодические.
\item
  \textbf{B.} Рекурсивные.
\item
  \textbf{C.} Состязательные.
\item
  \textbf{D.} Продолжающиеся.
\end{itemize}

\subsubsection{Вопрос 5: в чём состоит компромисс между исследованием и использованием знаний?}\label{q5-what-is-the-explorationexploitation-tradeoff}

\subsubsection{Вопрос 6: что такое политика?}\label{q6-what-is-a-policy}

\subsubsection{Вопрос 7: что такое методы на основе функции ценности?}\label{q7-what-are-value-based-methods}

\subsubsection{Вопрос 8: что такое методы на основе политики?}\label{q8-what-are-policy-based-methods}

\subsection{Ответы и объяснения}\label{answer-key-and-explanations}

\subsubsection{Вопрос 1: что такое обучение с подкреплением?}\label{q1-what-is-reinforcement-learning-1}

Обучение с подкреплением --- это \textbf{подход к решению задач управления (принятия решений)}, при котором создаются агенты, обучающиеся взаимодействовать со средой методом проб и ошибок и \textbf{получающие вознаграждения (положительные или отрицательные) в качестве единственной обратной связи}.

\subsubsection{Вопрос 2: опишите цикл RL}\label{q2-define-the-rl-loop-1}

\textbf{Ответ:} B.

\begin{itemize}
\item
  \textbf{Вариант A:} На каждом шаге агент получает от среды \textbf{состояние s0}. По \textbf{состоянию s0} он выбирает \textbf{действие a0}, например движение вправо. Среда переходит в \textbf{новое состояние s1} и выдаёт агенту \textbf{вознаграждение r1}.
\item
  \textbf{Вариант C:} На каждом шаге агент получает от среды \textbf{состояние s0}. По \textbf{состоянию s0} он выбирает \textbf{действие a0}, например движение вправо. Среда переходит в \textbf{новое состояние s1} и выдаёт агенту \textbf{вознаграждение r1}.
\end{itemize}

\subsubsection{Вопрос 3: чем состояние отличается от наблюдения?}\label{q3-whats-the-difference-between-a-state-and-an-observation-1}

\textbf{Ответ:} A, D, E, H.

\begin{itemize}
\item
  \textbf{Вариант E (верный):} Нам известна вся информация о шахматной доске.
\item
  \textbf{Вариант F:} Нам известна вся информация о шахматной доске.
\item
  \textbf{Вариант G:} Мы видим только часть уровня рядом с персонажем, поэтому получаем наблюдение.
\item
  \textbf{Вариант H (верный):} Мы видим только часть уровня рядом с персонажем.
\end{itemize}

\subsubsection{Вопрос 4: какие существуют два типа задач RL?}\label{q4-a-task-is-an-instance-of-a-reinforcement-learning-problem.-what-are-the-two-types-of-tasks-1}

\textbf{Ответ:} A, D.

\begin{itemize}
\item
  \textbf{Вариант A (верный):} Эпизодическая задача имеет начало и конец (терминальное состояние). Это образует эпизод --- последовательность состояний, действий, вознаграждений и следующих состояний. Например, в Super Mario Bros эпизод начинается с запуска нового уровня и завершается, когда персонаж погибает или доходит до конца.
\item
  \textbf{Вариант D (верный):} Продолжающиеся задачи идут бесконечно (терминального состояния нет). В таких задачах агент должен учиться выбирать оптимальные действия, одновременно продолжая взаимодействовать со средой.
\end{itemize}

\subsubsection{Вопрос 5: в чём состоит компромисс между исследованием и использованием знаний?}\label{q5-what-is-the-explorationexploitation-tradeoff-1}

В RL необходимо \textbf{соблюдать баланс между изучением среды и использованием уже известных сведений о ней}.

\begin{itemize}
\item
  \emph{Исследование (exploration)} --- изучение среды с помощью \textbf{случайных действий для получения дополнительной информации}.
\item
  \emph{Использование знаний (exploitation)} --- \textbf{применение известной информации для максимизации вознаграждения}.
\end{itemize}

\CourseFigure{expexpltradeoff.jpg}{Исследование и использование знаний: обобщающая схема}

\subsubsection{Вопрос 6: что такое политика?}\label{q6-what-is-a-policy-1}

\begin{itemize}
\tightlist
\item
  Политика $\pi$ --- \textbf{«мозг» агента}. Она определяет действие, которое следует выбрать в текущем состоянии, а значит, задаёт поведение агента в каждый момент времени.
\end{itemize}

\CourseFigure{policy_1.jpg}{Политика отображает текущее состояние в действие}

\subsubsection{Вопрос 7: что такое методы на основе функции ценности?}\label{q7-what-are-value-based-methods-1}

\begin{itemize}
\tightlist
\item
  Методы на основе функции ценности --- один из основных подходов к решению задач RL.
\item
  Вместо непосредственного обучения политики \textbf{мы обучаем функцию ценности, сопоставляющую состоянию ожидаемую ценность нахождения в нём}.
\end{itemize}

\subsubsection{Вопрос 8: что такое методы на основе политики?}\label{q8-what-are-policy-based-methods-1}

\begin{itemize}
\tightlist
\item
  В \emph{методах на основе политики} мы \textbf{непосредственно обучаем функцию политики}.
\item
  Такая функция \textbf{сопоставляет каждому состоянию наиболее подходящее действие} или \textbf{распределение вероятностей по доступным действиям}.
\end{itemize}

Поздравляем с прохождением теста! Если какие-то темы вызвали затруднения, ещё раз прочитайте соответствующие разделы для закрепления материала. Но \textbf{не переживайте}: в курсе мы неоднократно вернёмся к этим понятиям и \textbf{закрепим теорию практическими заданиями}.

\section{Заключение}\label{conclusion}

Поздравляем с завершением модуля! \textbf{Это был самый объёмный модуль}, и мы рассмотрели много новой информации. Поздравляем и с выполнением практической части: вы только что обучили первых агентов Deep RL и поделились ими с сообществом!

\textbf{Если некоторые понятия пока кажутся запутанными, это нормально}. Так было и у меня, и у всех, кто начинал изучать RL.

\textbf{Не торопитесь: прежде чем двигаться дальше, постарайтесь как следует освоить материал}. Важно иметь прочную базу перед переходом к более увлекательным задачам.

Конечно, в ходе курса мы будем снова использовать и объяснять эти термины, но лучше понять их до перехода к следующим модулям.

В следующем (бонусном) модуле мы закрепим пройденное, \textbf{обучив пса Huggy приносить палку}.

После этого вы сможете с ним поиграть.

\CourseFigure{huggy-bonus.jpg}{Пёс Huggy: обучающая среда из бонусного модуля}

Наконец, мы были бы рады \textbf{узнать ваше мнение о курсе и о том, как его можно улучшить}. Если хотите поделиться отзывом, пожалуйста, \href{https://forms.gle/BzKXWzLAGZESGNaE9}{заполните анкету}.

\subsubsection{Продолжайте учиться, у вас всё получится}\label{keep-learning-stay-awesome-1}

\section{Дополнительная литература}\label{additional-readings}

Эти материалы \textbf{необязательны}, но помогут глубже разобраться в теме.

\subsection{Глубокое обучение с подкреплением}\label{deep-reinforcement-learning}

\begin{itemize}
\tightlist
\item
  \href{http://incompleteideas.net/book/RLbook2020.pdf}{Reinforcement Learning: An Introduction, Richard Sutton и Andrew G. Barto, главы 1, 2 и 3}
\item
  \href{https://youtu.be/2GwBez0D20A}{Серия Foundations of Deep RL: лекция 1 о MDP, точных методах решения и RL с максимальной энтропией, Pieter Abbeel}
\item
  \href{https://spinningup.openai.com/en/latest/spinningup/rl_intro.html}{Spinning Up RL от OpenAI, часть 1: ключевые понятия RL}
\end{itemize}

\subsection{Gym}\label{gym}

\begin{itemize}
\tightlist
\item
  \href{https://blog.paperspace.com/getting-started-with-openai-gym/}{Начало работы с OpenAI Gym: основные строительные блоки}
\item
  \href{https://www.gymlibrary.dev/content/environment_creation/}{Создание собственной среды Gym}
\end{itemize}

\end{document}
